% Part: computability
% Chapter: recursive-functions
% Section: pr-relations

\documentclass[../../../include/open-logic-section]{subfiles}

\begin{document}

\olfileid[pt]{cmp}{rec}{prr}
\olsection{Relações recursivas primitivas}


\begin{defn}
Uma relação $R(\vec x)$ é dita recursiva primitiva se sua função
característica,
\[
\Char{R}(\vec x) = \left\{
  \begin{array}{ll}
  1 & \mbox{se $R(\vec x)$} \\
  0 & \mbox{caso contrário}
  \end{array}
\right.
\]
é recursiva primitiva.
\end{defn}

Em outras palavras, quando se fala de uma relação recursiva primitiva
$R(\vec x)$, faz-se referência a uma relação da forma
$\Char{R}(\vec x) = 1$, onde $\Char{R}$ é uma função recursiva primitiva que,
para qualquer entrada, devolve 1 ou 0. Por exemplo, a relação
$\fn{IsZero}(x)$, que vale se, e somente se, $x = 0$, corresponde à função
$\Char{\fn{IsZero}}$, definida por recursão primitiva mediante
\begin{align*}
\Char{\fn{IsZero}}(0) & = 1,\\
\Char{\fn{IsZero}}(x+1) & = 0.
\end{align*}

Deve estar claro que se podem compor relações com outras funções recursivas
primitivas. Portanto, também são recursivas primitivas:
\begin{enumerate}
\item A relação de igualdade, $x = y$, definida por
  $\fn{IsZero}(\left|x-y\right|)$.
\item A relação menor ou igual, $x \leq y$, definida por
  $\fn{IsZero}(x \tsub y)$.
\end{enumerate}


\begin{prop}
  O conjunto das relações recursivas primitivas é fechado sob operações
  booleanas; isto é, se $P(\vec x)$ e $Q(\vec x)$ são recursivas primitivas,
  então também o são
  \begin{enumerate}
  \item $\lnot P(\vec x)$
  \item $P(\vec x) \land Q(\vec x)$
  \item $P(\vec x) \lor Q(\vec x)$
  \item $P(\vec x) \lif Q(\vec x)$
  \end{enumerate}
\end{prop}

\begin{proof}
  Suponha que $P(\vec x)$ e $Q(\vec x)$ sejam recursivas primitivas, O.s.\@,
  que suas funções características $\Char{P}$ e $\Char{Q}$ o sejam. Precisamos
  mostrar que as funções características de $\lnot P(\vec x)$ etc. também são
  recursivas primitivas.
  \[
  \Char{\lnot P}(\vec x) = \begin{cases}
    0 & \text{se $\Char{P}(\vec x) = 1$}\\
    1 & \text{caso contrário}
  \end{cases}
  \]
  Podemos definir $\Char{\lnot P}(\vec x)$ como $1 \tsub \Char{P}(\vec x)$.
  \[
  \Char{P \land Q}(\vec x) = \begin{cases}
    1 & \text{se $\Char{P}(\vec x) = \Char{Q}(\vec x) = 1$}\\
    0 & \text{caso contrário}
  \end{cases}
  \]
  Podemos definir $\Char{P \land Q}(\vec x)$ como $\Char{P}(\vec x) \cdot
  \Char{Q}(\vec x)$ ou como $\fn{min}(\Char{P}(\vec x), \Char{Q}(\vec
  x))$. De modo semelhante,
  \begin{align*}
    \Char{P \lor Q}(\vec x) & = \fn{max}(\Char{P}(\vec x), \Char{Q}(\vec x)) \text{ e}\\
    \Char{P \lif Q}(\vec x) & = \fn{max}(1 \tsub
  \Char{P}(\vec x), \Char{Q}(\vec x)).
  \end{align*}
\end{proof}

\begin{prop}
  O conjunto das relações recursivas primitivas é fechado sob quantificação
  limitada; O.s.\@, se $R(\vec x, z)$ é uma relação recursiva primitiva,
  então também o são as relações
  \begin{align*}
    & \bforall{z < y}{R(\vec x, z)} \text{ e}\\
    & \bexists{z < y}{R(\vec x, z)}.
  \end{align*}
  $\bforall{z < y}{R(\vec x, z)}$ vale para $\vec x$ e $y$ se, e somente se,
  $R(\vec x, z)$ vale para todo~$z$ menor que~$y$, e analogamente para
  $\bexists{z < y}{R(\vec x, z)}$.
\end{prop}

\begin{proof}
  Por convenção, tomamos $\bforall{z < 0}{R(\vec x, z)}$ como verdadeiro
  (pela razão trivial de que não há $z$ menor que~$0$) e
  $\bexists{z < 0}{R(\vec x, z)}$ como falso. Um quantificador universal
  limitado funciona como um produto finito ou um mínimo iterado; O.s.\@, se
  $P(\vec x, y) \defiff \bforall{z < y}{R(\vec x, z)}$, então
  $\Char{P}(\vec x, y)$ pode ser definida por
  \begin{align*}
    \Char{P}(\vec x, 0) & = 1\\
    \Char{P}(\vec x, y+1) & =
    \fn{min}(\Char{P}(\vec x, y), \Char{R}(\vec x, y)).
  \end{align*}
  A quantificação existencial limitada pode ser definida analogamente usando
  $\fn{max}$. Alternativamente, pode ser definida a partir da quantificação
  universal limitada mediante a equivalência
  $\bexists{z < y}{R(\vec x, z)} \liff
  \lnot \bforall{z < y}{\lnot R(\vec x, z)}$. Observe, por exemplo, que um
  quantificador limitado da forma $\bexists{x \leq y}{\dots x\dots}$ é
  equivalente a $\bexists{x < y+1}{\dots x \dots}$.
\end{proof}

\begin{prob}
  Mostre que a relação ternária $x \equiv y \mod n$ (congruência módulo~$n$)
  é recursiva primitiva.
\end{prob}

Outra função recursiva primitiva útil é a função condicional
$\fn{cond}(x,y,z)$, definida por
\begin{align*}
  \fn{cond}(x,y,z) & = \begin{cases}
  y & \text{se $x = 0$} \\
  z & \text{caso contrário}.
\end{cases}
\intertext{Ela é definida recursivamente por}
\fn{cond}(0,y,z) & = y,\\
\fn{cond}(x+1,y,z) & = z.
\end{align*}
Podemos usá-la para justificar definições por casos de funções recursivas
primitivas a partir de relações recursivas primitivas:

\begin{prop}
Se $g_0(\vec x)$, \dots,~$g_m(\vec x)$ são funções recursivas primitivas e
$R_0(\vec x)$, \dots, $R_{m-1}(\vec x)$ são relações recursivas primitivas,
então a função $f$ definida por
\[
f(\vec x) = \begin{cases}
    g_0(\vec x) & \text{se $R_0(\vec{x})$} \\
    g_1(\vec x) & \text{se $R_1(\vec{x})$ e não $R_0(\vec{x})$} \\
    \vdots & \\
    g_{m-1}(\vec x) & \text{se $R_{m-1}(\vec{x})$ e nenhuma das anteriores
      vale}
    \\
    g_m(\vec x) & \mbox{caso contrário}
\end{cases}
\]
também é recursiva primitiva.
\end{prop}

\begin{proof}
  Quando $m = 1$, essa é simplesmente a função definida por
  \[
  f(\vec x) = \fn{cond}(\Char{\lnot R_0}(\vec x),g_0(\vec x),g_1(\vec
  x)).
  \]
  Para $m$ maior que $1$, basta compor definições dessa forma.
\end{proof}
\end{document}
